% =====================================================================
% Documento independiente: tabla de atención a las observaciones
% del evaluador del trabajo de grado. No forma parte del informe final.
% Requiere dos pasadas de compilación para que la tabla fije sus anchos.
% =====================================================================
\documentclass[11pt]{article}

\usepackage[utf8]{inputenc}
\usepackage[T1]{fontenc}
\usepackage{lmodern}
% es-noshorthands: necesario para que las comillas rectas (") de las
% transcripciones textuales del evaluador se compongan tal cual y no las
% capture el mecanismo de atajos de babel-spanish.
\usepackage[spanish,es-noshorthands]{babel}
\usepackage[landscape,letterpaper,top=1.8cm,bottom=2.2cm,left=1.6cm,right=1.6cm]{geometry}
\usepackage{longtable}
\usepackage{array}
\usepackage{ragged2e}
\usepackage{booktabs}
\usepackage{xcolor}
\usepackage{colortbl}
\usepackage[hidelinks]{hyperref}

% Columnas de párrafo con alineación a la izquierda y sin partición agresiva
\newcolumntype{C}[1]{>{\RaggedRight\arraybackslash}p{#1}}
\newcolumntype{Z}[1]{>{\centering\arraybackslash}p{#1}}

\definecolor{gristabla}{gray}{0.87}

\setlength{\LTcapwidth}{\textwidth}
\renewcommand{\arraystretch}{1.18}

\pagestyle{plain}

\newcommand{\titulotg}{Evolución arquitectónica y aseguramiento de calidad de la plataforma Paralegales mediante una arquitectura híbrida Firebase y AWS}
\newcommand{\autortg}{Juan Felipe Monsalve Vargas}
% Fecha fijada de forma explícita: el contenedor de compilación corre en UTC
% y \today adelantaría un día.
\newcommand{\fecha}{1 de agosto de 2026}

\begin{document}

\begin{center}
  {\Large\bfseries Tabla de atención a las observaciones del evaluador}\\[0.7em]
  {\normalsize\itshape \titulotg}\\[0.9em]
  {\normalsize \autortg}\\[0.3em]
  {\small Universidad del Valle, Escuela de Ingeniería de Sistemas y Computación, Sede (06) Tuluá}\\[0.3em]
  {\small \fecha}
\end{center}

\vspace{0.6em}

\noindent
Este documento relaciona cada una de las siete observaciones emitidas por el evaluador con la corrección
efectivamente aplicada al informe final y con el lugar donde puede verificarse. Los números de página
corresponden a la numeración impresa en el pie de página del informe. Las observaciones cuarta y quinta
se presentan en una sola fila, porque recaen sobre el mismo conjunto de figuras y se atendieron de forma
conjunta. Cuando una observación quedó atendida solo en parte, se declara como tal.

\vspace{0.8em}

\begingroup
\small
\begin{longtable}{@{}Z{1.0cm} C{6.7cm} C{9.2cm} C{6.2cm}@{}}
  \toprule
  \rowcolor{gristabla}
  \textbf{N.\textsuperscript{o}} &
  \textbf{Observación del evaluador} &
  \textbf{Corrección aplicada} &
  \textbf{Ubicación en el informe} \\
  \midrule
  \endfirsthead

  \multicolumn{4}{@{}l@{}}{\small\itshape (continuación de la página anterior)}\\[0.3em]
  \toprule
  \rowcolor{gristabla}
  \textbf{N.\textsuperscript{o}} &
  \textbf{Observación del evaluador} &
  \textbf{Corrección aplicada} &
  \textbf{Ubicación en el informe} \\
  \midrule
  \endhead

  \midrule
  \multicolumn{4}{@{}r@{}}{\small\itshape (continúa en la página siguiente)}\\
  \endfoot

  \bottomrule
  \endlastfoot

  % ---------------------------------------------------------------- 1
  1 &
  «EL RESUMEN ESTA MUY LARGO Y HAY MUCHAS PALABRAS CLAVE. AÑADAN EL ABSTRACT TAMBIEN» &
  Atendida. El resumen se condensó de cuatro párrafos a uno solo de 163 palabras, conservando el
  planteamiento del problema, la estrategia adoptada y el resultado obtenido. Las palabras clave se
  redujeron de ocho a cinco, dejando las que identifican los ejes del trabajo. Se añadió el
  \textit{Abstract} en inglés, con sus \textit{Keywords} correspondientes. &
  Página 1 (Resumen y «Palabras clave») y página 2 (\textit{Abstract} y \textit{Keywords}), en las
  páginas preliminares. \\
  \midrule

  % ---------------------------------------------------------------- 2
  2 &
  «en el documento se redacta mucho con guiones " - " eso no es una escritura natural para nosotros,
  y suele ser usado por la IA, revisen la escritura de todo el documento» &
  Atendida. Se revisó la redacción del documento completo, capítulo por capítulo, y se eliminó el uso
  de la raya para intercalar aclaraciones, que era el hábito que el evaluador señaló. Los pasajes
  afectados, alrededor de cincuenta y cinco, se reescribieron variando el recurso según lo que pedía
  cada frase: comas cuando la aclaración era breve, paréntesis cuando convenía apartarla del hilo
  principal, dos puntos cuando introducía una enumeración y, en varios casos, la división de la oración
  en dos. Se conservaron únicamente los guiones que corresponden a palabras compuestas y al nombre de
  un patrón de arquitectura. Hoy no queda en el informe ningún inciso construido de ese modo. &
  Corrección transversal a todo el documento, apreciable en cualquiera de sus capítulos. \\
  \midrule

  % ---------------------------------------------------------------- 3
  3 &
  «MARCO TEORICO DE UNA SOLA SECCION?» &
  Atendida. El marco teórico pasó de una a tres subsecciones: «Teoría General de Sistemas»
  (preexistente), «Leyes de la evolución del software» (Lehman) y «Teoría de la deuda técnica»
  (Cunningham; Kruchten, Nord y Ozkaya), con las referencias correspondientes añadidas a la
  bibliografía. Cada subsección nueva cierra con su aplicación al caso Paralegales. Además, ese marco
  se invoca posteriormente en el capítulo de resultados (objetivo B) y en las conclusiones, de modo
  que gobierna el resto del documento y no queda aislado. &
  Subsección 5.2 «Marco Teórico», páginas 26 a 29: «Leyes de la evolución del software» en la página
  27 y «Teoría de la deuda técnica» en la página 28. Invocación posterior en el objetivo B, páginas 90
  y 91 («partidas de deuda de mayor interés», «deuda estructuralmente invisible»), y en las
  conclusiones, página 100 («programa de tipo E», «pago selectivo»). \\
  \midrule

  % ---------------------------------------------------------------- 4 y 5
  4 y 5 &
  «LOS PANTALLAZOS SE PUEDEN MEJORAR, REVISARLOS TODOS»

  \vspace{0.5em}
  «en graficas como esta, o la 33, intuyo que se muestran datos con usuarios reales. Tener cuidado con
  los datos personales de la gente, sobre todo por cuidar a la empresa, revisen esto» &
  Ambas observaciones se atendieron de forma conjunta, al recaer sobre el mismo conjunto de figuras. Se
  revisaron una por una todas las capturas del informe, valorando en cada caso tanto lo que aportaban
  al argumento como la información que dejaban a la vista. De esa revisión resultó la intervención de
  diez figuras, en las que se cubrieron los datos que permitían identificar a personas concretas:
  nombres de partes de procesos judiciales, correos de terceros, alias y nombres de miembros del
  equipo, avatares con rostro, identificadores de usuario y de pago, y credenciales de acceso a
  servicios. Se conservaron, en cambio, los datos de registro público, como los números de radicado,
  las entidades y los juzgados, por ser parte del objeto del sistema y no información personal
  sensible. Queda pendiente extender ese mismo tratamiento a la Figura 3 y suprimir el identificador
  de cuenta y el principal de acceso que aún se aprecian en las Figuras 6, 27 y 28, así como los
  correos del autor en las Figuras 10 y 13. &
  Figuras intervenidas: apreciable en las figuras del informe, comparadas con la versión entregada
  previamente.

  \vspace{0.5em}
  Pendiente: Figura 3 (página 43), Figura 6 (página 50), Figura 10 (página 54), Figura 13 (página 57),
  Figura 27 (página 73) y Figura 28 (página 74). \\
  \midrule

  % ---------------------------------------------------------------- 6
  6 &
  «segun entiendo, la automatización se concentra solo en las reglas de seguridad de Firestore. Para
  un trabajo cuyo título incluye "aseguramiento de calidad" como eje central considero que se podria
  mejorar» &
  Atendida. En el marco conceptual se definió «aseguramiento de calidad» como proceso, según la norma
  IEEE 730 y el modelo ISO/IEC 25010, y no como una batería de pruebas. El capítulo se reencuadró en
  ocho mecanismos distribuidos a lo largo del ciclo de vida, acompañados de una tabla matriz que
  registra, para cada uno, el momento de ejecución, el grado de automatización y la evidencia; cuatro
  de esos mecanismos operan como compuertas automáticas. Se corrigió la afirmación de que el
  \textit{backend} carecía de pruebas, atribuyendo la línea base en Go al equipo de la empresa. Las
  pruebas \textit{end-to-end} pasaron de declararse «no consolidadas» a «implementadas, no
  integradas». La observabilidad se incorporó como control de calidad posdespliegue. Las limitaciones
  y el trabajo futuro se reescribieron en tres órdenes. &
  Definición con IEEE 730 e ISO/IEC 25010 en el marco conceptual, subsección 5.1.11, página 25.
  Subsección 9.1 «Enfoque y Estrategia de Aseguramiento de Calidad», página 77. Tabla 4, «Mecanismos
  de aseguramiento de calidad, momento de ejecución y grado de automatización», página 85. Pruebas
  \textit{end-to-end} en la subsección 9.4, página 82. Limitaciones y trabajo futuro en las páginas 85
  y 86. \\
  \midrule

  % ---------------------------------------------------------------- 7
  7 &
  «Se menciona SRE y menciona explícitamente MTTR y error budgets en el marco teórico, porque no se
  reporta ningún dato con estos indicadores?» &
  Atendida. En el marco conceptual se delimitó qué se toma de la práctica de SRE (instrumentación,
  indicadores de nivel de servicio y alerta por umbral) y qué no se toma (SLOs, \textit{error budgets}
  y MTTR). En el capítulo de resultados se añadió una tabla con los cinco indicadores efectivamente
  implementados (\texttt{RJURequestCount}, \texttt{RJU5xxCount}, \texttt{RJU4xxCount},
  \texttt{RJULatencyMs} y \texttt{RJUStatusDistribution}) y las dos compuertas de alerta con sus
  umbrales y ventanas, verificados contra la plantilla de CloudFormation. La ausencia de SLOs,
  \textit{error budgets} y MTTR se declara como limitación, indicando su causa, y como trabajo futuro,
  con eco en las conclusiones. &
  Delimitación de lo tomado de SRE en el marco conceptual, subsección 5.1.9, página 24. Tabla 5,
  «Indicadores de nivel de servicio implementados sobre las integraciones externas, umbrales
  configurados y acciones asociadas», páginas 92 y 93. Declaración de la limitación y del trabajo
  futuro en la página 94. Eco en las conclusiones, página 99. \\

\end{longtable}
\endgroup

\vspace{0.4em}

\noindent
\small
\textbf{Nota sobre la numeración.} Las páginas citadas son las impresas en el pie de página del
informe. Si se navega el documento digital, cuyo conteo empieza en la portada, cada una de esas páginas
equivale a la posición dos unidades mayor: por ejemplo, la página impresa 92 se abre en la posición 94.

\end{document}
